\documentclass[a4paper,10pt]{article}
\usepackage[spanish]{babel}
\usepackage{
			ucs
			,amsfonts
			,amssymb
			,acronym
			,verbatim
			,rotating
			,graphicx
			,enumerate
			,fancyhdr
            ,mathrsfs
            ,array
			}

\usepackage{tikz}
\usetikzlibrary{svg.path}

\pagestyle{fancy}
\addtolength{\textheight}{1.0in}
\addtolength{\textwidth}{0.5in}
\addtolength{\voffset}{-1.0in}

\newcommand{\cat}[1][]{%
    \tikz \fill [scale=1ex/30,yscale=-1,#1] svg "M495.959,867.638c-3.318-4-1.121-7.686,7.098-11.904c7.322-3.76,8.223-4.871,11.365-14.049
		c5.646-16.496,9.951-21.084,24.289-25.893c4.883-1.637,7.111-3.217,6.732-4.771c-3.148-12.889-1.992-16.166,4.48-12.701
		c2.174,1.162,3.924,0.754,6.943-1.619c4.547-3.576,5.105-3.137,6.68,5.264c0.725,3.859,0.445,6.785-0.791,8.273
		c-1.186,1.43-1.496,4.199-0.816,7.295c0.895,4.07,0.186,6.26-3.836,11.861c-6.01,8.367-6.746,18.059-1.752,23.053
		c3.975,3.973,4.066,6.363,0.248,6.363c-3.869,0-8.201-4.926-8.201-9.322c0-4.723-4.635-6.33-7.391-2.564
		c-2.844,3.889-2.738,7.115,0.23,7.115c1.312,0,2.387,1.074,2.387,2.387c0,1.92-2.465,2.385-12.635,2.385
		c-6.947,0-13.047-0.668-13.551-1.484c-0.541-0.875-3.803-0.502-7.951,0.912c-7.744,2.639-11.631,5.996-8.896,7.686
		c0.916,0.566,1.668,1.703,1.668,2.527C502.26,870.757,498.104,870.222,495.959,867.638z";~%
}
\renewcommand{\cat}{\maltese}




\fancyhf{}
\fancyhead[L]{{\bfseries Programaci\'on I Curso 2012. Pr\'actica 2 \hfill Pseudo-C\'odigo - Resoluci\'on}}
\fancyfoot[LE,RO]{\thepage}

\fancypagestyle{firstpage}{
\fancyhf{}

\fancyhead[L]{\includegraphics[width=1.75cm]{ucalp.jpg}  \\
\small{UCALP. Facultad de Ciencias Exactas e Ingenier\'ia} \\
{\bfseries Programaci\'on I Curso 2012. Pr\'actica 2 \hfill Pseudo-C\'odigo - Resoluci\'on}}
\fancyfoot[LE,RO]{\thepage}
}

\thispagestyle{firstpage}

\begin{document}

~\\ \\ \\ \\
\begin{itemize}

\item[3]
    \begin{itemize}
        \item[d-] \begin{verbatim} Gamma := tar;\end{verbatim} INV\'ALIDA. No se puede asignar un valor a una constante.
        \item[e-] \begin{verbatim} rango := respuesta;\end{verbatim} INV\'ALIDA. Los tipos de las variables no concuerdan.
        \item[g-] \begin{verbatim} tar := mx9;\end{verbatim} V\'ALIDA. El valor de \verb=mx9= se convierte autom\'aticamente a real (al reves no se puede).  
        \item[h-] \begin{verbatim} simbolo := `t`;\end{verbatim} V\'ALIDA. Los tipos son correctos y se puede realizar la asignaci\'on.
    \end{itemize}

\item[4.b-] ~~
    \begin{table}[!h]
      \begin{tabular}{|c||c|c|c|}
        \hline
        Nro. de Linea & \verb=x= & \verb=y= & \verb=z= \\
        \hline \hline
        1 & 27 & -- & -- \\
        2 & 27 & 343 & -- \\
        3 & 27 & 343 & 70 \\
        4 & 30 & 343 & 70 \\
        5 & 30 & 100 & 70 \\
        6 & 200 & 100 & 70 \\
        7 & 200 & 200 & 70 \\
        \hline
      \end{tabular}
    \end{table}

\item[5]
    \begin{itemize}
        \item[a-] \begin{verbatim} x1 := b * b / -7 / 1.45;\end{verbatim} V\'ALIDA. Se asigna a \verb=x= un valor en funci\'on del valor de \verb=b=
        \item[b-] \begin{verbatim} x := sqrt( abs( 5 * 5 - 4 * 4));\end{verbatim} V\'ALIDA. Se asigna a \verb=x= en valor 3.
        \item[c-] \begin{verbatim} -v := x * y - total;\end{verbatim} ERROR. \verb=-v= no es un nombre de variable v\'alido suponiendo que \verb=v= es la variable, no se pueden aplicar funciones o operadores unarios del lado izquierdo de una sentencia de asignaci\'on.
        \item[d-] \begin{verbatim} precioTotal - descuento := precioPagar;\end{verbatim} ERROR. 
    \end{itemize}	

\item[7.c-]
\begin{small}
\begin{verbatim}
program siete_c;

var
  i,j:integer;

begin
   for i:= 0 to 4 do
   begin
     for j:= 0 to (3 - (i mod 4)) do
         write(' ');
     for j:= 0 to (2 * (i mod 4)) do
         write('*');
     for j:= 0 to (3 - (i mod 4)) do
         write(' ');
     writeln();
    end;
   readln;
end.
\end{verbatim}
\end{small}

\item[9]
\begin{small}
\begin{verbatim}
program nueve;

var
  total:integer;

begin
   write('Ingrese la cant. de dinero: $');
   readln(total);
   writeln('Billtes de $100: ', total div 100);
   total := total mod 100;
   writeln('Billtes de $50: ', total div 50);
   total := total mod 50;
   writeln('Billtes de $20: ', total div 20);
   total := total mod 20;
   writeln('Billtes de $10: ', total div 10);
   total := total mod 10;
   writeln('Billtes de $5: ', total div 5);
   total := total mod 5;
   writeln('Billtes de $2: ', total div 2);
   total := total mod 2;
   writeln('Monedas de $1: ', total);

   readln;
end.
\end{verbatim}
\end{small}

\item[10]
\begin{small}
\begin{verbatim}
program diez;

var
  cuenta,:integer;

begin
   write('Ingrese la cant. de segundos: ');
   readln(cuenta);

   writeln((cuenta div 60) div 60,' : ',cuenta div 60 mod 60,' : ',((cuenta mod 60) mod 60));

   readln;
end.
\end{verbatim}
\end{small}
\item[15]
\begin{small}
\begin{verbatim}
program quince;

var
  ret,ab,i:integer;

begin
   write('Ingrese un numero: ');
   readln(ab);

   for i:= 1 to trunc(ab/2) do
      if (ab mod i = 0) then
         ret:= ret + i;

   if (ret > ab) then
      writeln('El numero es abundante')
   else
      writeln('El numero NO es abundante');

   readln;
end.
\end{verbatim}
\end{small}
\item[17]
\begin{small}
\begin{verbatim}
program d_siete;

var
  ret,def,i:integer;

begin
   write('Ingrese un numero: ');
   readln(def);

   for i:= 1 to trunc(def/2) do
      if (def mod i = 0) then
         ret:= ret + i;

   if (ret < def) then
      writeln('El numero es defectivo')
   else
      writeln('El numero NO es defectivo');

   readln;
end.
\end{verbatim}
\end{small}
\item[19]
    \begin{itemize}
      \item[3]
\begin{small}
\begin{verbatim}
program dn_tres;

var
  b,h:real;

begin
   write('Ingrese la altura del rectangulo: ');
   readln(h);
   write('Ingrese la base del rectangulo: ');
   readln(b);

   writeln('El area del rectangulo es: ',b*h);

   readln;
end.
\end{verbatim}
\end{small}
      \item[4]
\begin{small}
\begin{verbatim}
program dn_cuatro;

var
  b,h:real;

begin
   write('Ingrese la altura del triangulo: ');
   readln(h);
   write('Ingrese la base del triangulo: ');
   readln(b);

   writeln('El area del triangulo es: ',(b*h)/2);

   readln;
end.
\end{verbatim}
\end{small}
      \item[10]
\begin{small}
\begin{verbatim}
program dn_diez;

var
  km,pl,pg,h,m:real;

begin
   write('Ingrese el total de kilometros recorridos: ');
   readln(km);
   write('Ingrese el precio del litro de nafta: ');
   readln(pl);
   write('Ingrese el dinero gastado en nafta: ');
   readln(pg);
   write('Ingrese el tiempo total del viaje: ');
   readln(h,m);

   {consumo de gasolina cada 100km (en litros y pesos}
   writeln('Consumo de gasolina cada 100km: ',((pg/pl)*100)/km,' litros');
   writeln('Consumo de gasolina cada 100km: $',(((pg/pl)*100)/km)*pg,' pesos');
   {consumo de gasolina en litros y pesos cada km}
   writeln('Consumo de gasolina cada por km: ',(pg/pl)/km,' litros');
   writeln('Consumo de gasolina cada por km: $',((pg/pl)/km)*pg,' pesos');
   {velocidad media en km/h y m/s}
   writeln('Velocidad promedio: ',km/(h + m/60),' km/h');
   writeln('velocidad promedio: ',(km * 1000)/((h*60 + m)*60),' m/s');

   readln;
end.
\end{verbatim}
\end{small}
    \end{itemize}

\end{itemize}
\end{document}
